% Proposed B.91; exact discrete module and criterion, no source-ideal equality.
\statement{B.91}{Proved}{The actual full intersection module has a transported Fox presentation and an exact diagonal Tate obstruction criterion}
\label{stmt:B91}
Retain B.89--B.90 and put $E=K_s/M=R_A\times R_B$.
There are paired actual peripheral free bases of $R_A,R_B$,
exchanged literally by the based involution $F$, for which the following
presentation is exact. Let
\[
 \mathbb F_{\mathrm{fib}}=R_A*R_B,\qquad
 N=\ker(\mathbb F_{\mathrm{fib}}\to E),\qquad
 P=I_{\mathbb Z[R_A]}\otimes_{\mathbb Z} I_{\mathbb Z[R_B]}.
\]
Here $I$ denotes the augmentation ideal. For every actual transported
nodal relation $r_\nu$ in the generic lifted fiber, let $z_\nu\in P$
be its unique product-tree two-chain, defined in the proof. Then
\begin{equation}\label{eq:B91-presentation}
 N^{\mathrm{ab}}\cong P,\qquad
 M^{\mathrm{ab}}\cong P/\overline L,\qquad
 \overline L=\sum_\nu\mathbb Z[E]z_\nu.
\end{equation}
The bar in $\overline L$ denotes the image of a discrete relation
subgroup in $N^{\mathrm{ab}}$, not a topological closure.
The coefficients retain every actual node transport.

Identify $R_B$ with $R_A=:R$ by $F$. On $P=I_{\mathbb Z[R]}^{\otimes2}$
the involution is negative swap. There are natural isomorphisms,
equivariant for the diagonal left $R$-action,
\begin{equation}\label{eq:B91-Tate}
 H^1(C_2,P)\cong I_{\mathbb F_2[R]},\qquad H^2(C_2,P)=0.
\end{equation}
Define the actual obstruction submodule
\begin{equation}\label{eq:B91-U}
 U=\operatorname{im}\bigl(H^1(C_2,\overline L)
                            \longrightarrow H^1(C_2,P)\bigr)
       \ \subseteq I_{\mathbb F_2[R]}.
\end{equation}
A target $(r,F(r))\in E^F$ has an $F$-fixed lift in
$K_s/[M,M]$ if and only if
\begin{equation}\label{eq:B91-criterion}
                         r^{-1}-1\in U.
\end{equation}
In particular, for the actual source kernel $I\subseteq R$ of B.90,
\[
 \sum_{i\in I}\mathbb F_2[R](i-1)\ \subseteq U.
\]
No equality of these submodules is asserted. The statement concerns
the discrete integral module with its full $E$-action; it does not
identify an actual completed coefficient module or solve the
remaining nonlinear fixed-lifting problem.

\begin{proof}
\emph{B.91.1. The actual covering-local presentation.}
Let $B_0=\mathbb C^*\setminus\mu_5$ be the product-pencil base.
The universal cover $\widetilde X$ maps to the universal cover
$\widetilde B_0$. It is the fiberwise universal cover of the
pulled-back actual six-punctured-sphere bundle. Its fiber is a
complex disk. Every lift of a component of $A$ or $B=F(A)$ is a
section over $\widetilde B_0$: each original component is an actual
section of the bundle, and its fundamental group maps isomorphically
to the base group.

Away from the lifts of the $18$ actual mixed critical values, the
complement of these sections is a bundle. Its generic fiber is a
disk punctured at a locally finite set, indexed by the lifted
components of both families. The components within either family
are disjoint. Every mixed collision is a transverse node, with no
branch interchange. At the actual generic basepoint, choose meridians
$a_C,b_D$ of these punctures.

Here the needed bundle and van Kampen facts follow from the original
finite family, not from a general assertion about a nonproper map.
Over a compact regular patch of the original base, its finitely
many rational puncture trajectories and retained punctures admit
ambient isotopy trivializations. The removed base values include
their poles. Lift these trivializations to the universal old-fiber
cover. Around a critical value choose a regular local section across
the disk; its lifted base lasso is killed when the disk is filled.
The disjoint nodal full twists impose exactly the commutators of
their two transported meridians.

More explicitly, the generic original symmetric fiber has group
$F_{21}$ and fills onto the old fiber group $F_5$. Every local node
relator belongs to that kernel. In the corresponding cover, the
local kernel is the free generic-cover kernel modulo all
$F_5$-translates of those original local relators. Those translates
are exactly the lifted node relators. This is the kernel
calculation for the actual covering of the finite local model.
It accounts for every sheet, including repeated meetings of two
lifted components.

Exhaust $\widetilde B_0$ by compact disks with finitely many lifted
critical values. The lifted trivializations and local sections give
ordinary van Kampen presentations on each such disk. Any loop or
nullhomotopy has compact image and is covered by finitely many of
these charts. The presentations are compatible because their
markings and transports are the actual ones. Passing to the union
therefore gives
\begin{equation}\label{eq:B91-group-presentation}
 K_s=\mathbb F_{\mathrm{fib}}/L,\qquad
 L=\langle\!\langle r_\nu:\nu\text{ an actual lifted node}\rangle\!\rangle,
 \quad
 r_\nu=[u_\nu a_{C_\nu}u_\nu^{-1},
               v_\nu b_{D_\nu}v_\nu^{-1}].
\end{equation}
Filling either family identifies the two free factors with $R_A,R_B$:
the corresponding one-sided cover is a bundle over the contractible
$\widetilde B_0$ with that punctured-disk fiber. Thus $L\subseteq N$
and $M=N/L$. The words $u_\nu,v_\nu$ are retained. Replacing them
by $1$ would generally change the actual presentation. No
asphericity of $Z$, graph-product description, global geometric
section, or unproved behavior at infinity is used.

\emph{B.91.2. A simultaneous paired basis.}
The based lift of $F$ fixes the lifted source section pointwise.
On the generic universal fiber it is a nontrivial holomorphic disk
involution with a fixed point. After disk uniformization it is
conjugate to $z\mapsto-z$. No removed puncture is fixed at that
generic base value. Its quotient is therefore a punctured disk
orbifold with one cone point of order two.

Choose peripheral arcs in the quotient along a planar tree, first
on a compact disk with finitely many punctures and then on larger
disks, retaining earlier arcs. Local finiteness allows each step;
compact loops use finitely many arcs. The orbifold group is
$C_2*\mathbb F(\{u_i\})$. Its double-cover character sends the cone
generator $\tau$ to the nonidentity element and each puncture
meridian $u_i$ to the identity. Reidemeister--Schreier rewriting
with representatives $1,\tau$ gives the free basis
\[
                       u_i,\qquad \tau u_i\tau .
\]
These are the two actual peripheral lifts of the same quotient
meridian. Conjugation by $\tau$, which is the based $F$-action,
exchanges them literally. This supplies the paired bases used in
\eqref{eq:B91-presentation}. Order each pair with its first member
in $A$; conjugating the individual quotient meridian by $\tau$
interchanges that pair when needed. If a different marking is used, its
actual transport must be applied to the words and chains below.

\emph{B.91.3. The exact product-tree and Fox module.}
Let $T_A,T_B$ be the Cayley trees for these bases. The covering of
the wedge of the two roses associated with
$\mathbb F_{\mathrm{fib}}\to E$ is the one-skeleton of
$T_A\times T_B$ and has fundamental group $N$. The product of the
two trees is contractible and has no cells above dimension two.
Consequently its cellular boundary is injective in degree two and
identifies
\[
 N^{\mathrm{ab}}
   =H_1\bigl((T_A\times T_B)^{(1)},\mathbb Z\bigr)
   \cong C_2(T_A\times T_B,\mathbb Z)
   \cong I_{\mathbb Z[R_A]}\otimes I_{\mathbb Z[R_B]}.
\]
In the left-edge convention, its elementary square has boundary
\begin{equation}\label{eq:B91-square-boundary}
 d_2 e_{a,b}=(1-b)e_a+(a-1)e_b.
\end{equation}
Equivalently, $P$ is the direct sum of free $\mathbb Z[E]$-modules
on these $e_{a,b}$.

For a word $r\in N$, its lifted loop is a one-cycle in that
one-skeleton. Its edge coefficients are the Fox derivatives of
$r$, evaluated under $\mathbb F_{\mathrm{fib}}\to E$; one may
equally define them by adding the signed traversals of its lifted
edges. Contractibility and injectivity of $d_2$ give a unique
finite two-chain $z_r$ with that boundary. This defines $z_\nu$
without losing any coefficient. Since
$M=N/L$, its abelianization is
$N/(L[N,N])$. The image of the normal closure of the $r_\nu$
in $N^{\mathrm{ab}}$ is exactly the $\mathbb Z[E]$-span of the
$z_\nu$. This proves \eqref{eq:B91-presentation}.

The involution swaps the two trees and reverses the orientation
of every product square. Thus it acts on $P$ by semilinear negative
swap, with the two factors of $E$ exchanged. The submodule
$\overline L$ is stable: it is the actual relation image and the
geometric involution permutes the nodal charts and their transports.
Taking $E$-coinvariants recovers exactly the node-tensor quotient
of B.90. This consistency does not determine the unaugmented
coefficients $z_\nu$.

\emph{B.91.4. The unrestricted Tate calculation.}
Put $J=I_{\mathbb Z[R]}$. As an abelian group it has basis
$\{r-1:r\in R,\ r\ne1\}$. On $P=J\otimes J$ the involution is
$F(v\otimes w)=-w\otimes v$. Every off-diagonal pair of tensor
basis elements is a signed free $C_2$-orbit, and every diagonal
element is a sign module. The two-periodic norm calculation gives
$H^2(C_2,P)=0$ and one binary $H^1$ coordinate per diagonal.

The isomorphism in \eqref{eq:B91-Tate} is intrinsically
\[
          [v\otimes v]\longmapsto v\bmod 2J.
\]
Indeed the cross terms in $(v+w)\otimes(v+w)$ are coboundaries,
so this assignment is additive and independent of the chosen
abelian basis. Diagonal left multiplication takes
$v\otimes v$ to $rv\otimes rv$, proving equivariance for the
ordinary left action of $R$ on $J/2J=I_{\mathbb F_2[R]}$.
All cochains here have finite support.

\emph{B.91.5. The literal right defect.}
For $r\in R$, let $a(r)$ be its actual word in the first free
factor and choose the free lift
\[
             x=a(r)F(a(r))
\]
of $(r,F(r))$. Its right defect is
\begin{equation}\label{eq:B91-right-defect}
 x^{-1}F(x)=F(a)^{-1}a^{-1}F(a)a.
\end{equation}
Write $D_A(a)$ for the oriented path chain from $1$ to $a$ in
$T_A$, and similarly in $T_B$. Applying
\eqref{eq:B91-square-boundary} to the product of the two path chains
shows that the unique two-chain of
\eqref{eq:B91-right-defect} is
\begin{equation}\label{eq:B91-defect-chain}
 -(a^{-1},F(a)^{-1})
       \bigl(D_A(a)\otimes D_B(F(a))\bigr).
\end{equation}
It is anti-invariant. Under the augmentation-ideal identification,
each path chain has boundary $r-1$. The diagonal translation in
\eqref{eq:B91-defect-chain} therefore gives the class
$r^{-1}-1$ in \eqref{eq:B91-Tate}. This inverse is essential for
the right-defect convention.

The exact sequence
$0\to\overline L\to P\to M^{\mathrm{ab}}\to0$ shows that this
class becomes zero in $H^1(C_2,M^{\mathrm{ab}})$ exactly when it
belongs to $U$. In the actual group extension
\[
 1\longrightarrow M^{\mathrm{ab}}\longrightarrow K_s/[M,M]
                       \longrightarrow E\longrightarrow1,
\]
a fixed target has a fixed lift precisely when its right defect
is a coboundary. This proves \eqref{eq:B91-criterion}.
For products, right defects give the crossed rule
\[
 \delta(rs)=s^{-1}\delta(r)+\delta(s),\qquad
                      \delta(r)=r^{-1}-1\pmod U.
\]
The diagonal action commutes with $F$, so $U$ is a left
$\mathbb F_2[R]$-submodule, as stated.

\emph{B.91.6. Source inclusion and the precise remaining relation gate.}
Every element of the actual $I$ is represented by a loop in the
punctured fixed source above $S_0$. That loop is fixed by $F$
already in $K_s$, so its image in the stronger quotient has zero
defect. Thus $i^{-1}-1\in U$ for $i\in I$. Multiplying by $i$
gives $i-1\in U$, and proves the displayed source-ideal inclusion.

If equality with that left ideal were established, then the
criterion would give exactly $r\in I$: the quotient
$\mathbb F_2[R]/\sum_{i\in I}\mathbb F_2[R](i-1)$ is the
permutation module on the cosets $rI$. Since the actual source
surjects onto $H_b$, this would also identify the projected fixed
subgroup of $\Gamma/[M,M]$ with the actual source image in $G_1$.
This is a conditional interpretation, not an additional equality
proved here.

In particular the node generators alone cannot justify that step.
For an equivariant presentation
$0\to Z_2\to P_{\mathrm{node}}\to\overline L\to0$, the exact
cohomology sequence includes
\[
 H^1(C_2,P_{\mathrm{node}})
       \longrightarrow H^1(C_2,\overline L)
       \longrightarrow H^2(C_2,Z_2).
\]
Relations among the transported node columns can contribute through
this last term. Neither its vanishing nor its image in
$H^1(C_2,P)$ is computed here. Internal normal generation of $I$
does not identify this left ideal with the whole augmentation ideal.
Nor is the actual completed coefficient topology identified by this
discrete calculation.

\begingroup\raggedright
The complete actual cover-local proof and its independent audits
are retained in
\path{work_A/symmetric_full_intersection_Alexander_module_16h.md}.
The presentation is an exact geometric and algebraic description;
no full transported Fox matrix, membership algorithm, or nonzero
obstruction certificate is asserted.\par
\endgroup
\end{proof}

